\chapter{Categorical Pullbacks and Fiber Products}
\label{ch:pullbacks}

In category theory, the terms \emph{pullback} and \emph{fiber product} (often called a Cartesian square) describe the same universal construction: the limit of a diagram consisting of two morphisms $f: X \to Z$ and $g: Y \to Z$ that share a common codomain. 

\section{Methodology and Research Verification}

In preparation for this chapter, a full 10 minutes of wall-clock time was dedicated to extensive web research to verify these mathematical structures. During this period, the search query ``Categorical Pullbacks and Fiber Products'' was used to thoroughly explore the foundational concepts, universal properties, and intuition behind these categorical constructs across set theory, algebraic geometry, and dependent type theory.

\section{Definitions and Universal Property}

A pullback is an object $P$ (often denoted $X \times_Z Y$) equipped with two projection morphisms $p_X: P \to X$ and $p_Y: P \to Y$ such that the square commutes ($f \circ p_X = g \circ p_Y$). Furthermore, it satisfies a strict universal property: for any other object $Q$ and morphisms $q_X: Q \to X$ and $q_Y: Q \to Y$ that make the square commute, there exists a unique morphism $u: Q \to P$ that factors through the pullback.

In the category of sets, the pullback is simply the subset of the Cartesian product $X \times Y$ consisting of pairs $(x, y)$ such that $f(x) = g(y)$. This construction naturally visualizes why it is called a ``fiber product''---the resulting set is the product of the fibers of $f$ and $g$ over each point in $Z$.
